% Elegant ACM Conference Paper Template
% Clean and professional layout inspired by modern academic papers
% Based on ACM sigconf format with optimized typography

\documentclass[sigconf, nonacm, anonymous]{acmart}

%% Core packages for professional typography
\usepackage{microtype}      % Improved typography and spacing
\usepackage{graphicx}       % Enhanced graphics support
\usepackage{booktabs}       % Professional quality tables
\usepackage{amsmath}        % Advanced mathematical typesetting
\usepackage{algorithm}      % Algorithm environment
\usepackage{algpseudocode}  % Algorithm pseudocode formatting

%% Optional packages (uncomment as needed)
% \usepackage{subcaption}   % For subfigures
% \usepackage{xcolor}       % For colored text
% \usepackage{listings}     % For code listings
% \usepackage{hyperref}     % For hyperlinks (usually loaded by acmart)

%% Paper title - use title case
\title{Your Paper Title: A Clean and Professional Approach}

%% Author information (for anonymous submission)
\author{Anonymous Author(s)}
\affiliation{%
  \institution{Anonymous Institution}
  \city{Anonymous City}
  \country{Anonymous Country}
}
\email{anonymous@email.com}

%% For camera-ready version, replace with actual authors:
% \author{First Author}
% \affiliation{%
%   \institution{Your University}
%   \city{City Name}
%   \country{Country}
% }
% \email{first.author@university.edu}

%% Abstract - Keep it concise and informative (150-250 words)
\begin{abstract}
Write your abstract here. A good abstract should:
(1) clearly state the problem you're addressing,
(2) briefly describe your approach or methodology,
(3) highlight the main results or contributions, and
(4) indicate the significance of your work.

Keep the language clear and accessible, avoiding unnecessary jargon. The abstract is often the first (and sometimes only) part of your paper that readers will see, so make it count.
\end{abstract}

%% CCS Concepts (optional - for ACM publications)
% \begin{CCSXML}
% <ccs2012>
%   <concept>
%     <concept_id>10010147.10010178</concept_id>
%     <concept_desc>Computing methodologies~Machine learning</concept_desc>
%     <concept_significance>500</concept_significance>
%   </concept>
% </ccs2012>
% \end{CCSXML}
% \ccsdesc[500]{Computing methodologies~Machine learning}

%% Keywords (optional but recommended)
\keywords{keyword1, keyword2, keyword3, keyword4}

%% Start document
\begin{document}

\maketitle

%% Main content sections
\section{Introduction}

Your introduction goes here. Introduce the problem, motivation, and your contributions.

\subsection{Contributions}

List your main contributions:
\begin{itemize}
  \item First contribution
  \item Second contribution
  \item Third contribution
\end{itemize}

\section{Related Work}

Discuss related work and how your approach differs.

\section{Methodology}

Describe your approach and methodology.

\subsection{System Design}

Details about your system design.

\section{Evaluation}

Present your experimental results.

\subsection{Experimental Setup}

Describe your experimental setup.

\subsection{Results}

Present and discuss your results.

\begin{figure}[t]
  \centering
  \includegraphics[width=\columnwidth]{example-figure.pdf}
  \caption{Example figure caption. Describe what the figure shows.}
  \label{fig:example}
\end{figure}

\begin{table}[t]
  \caption{Example table caption}
  \label{tab:example}
  \centering
  \begin{tabular}{lcc}
    \toprule
    Method & Metric 1 & Metric 2 \\
    \midrule
    Baseline & 0.85 & 0.72 \\
    Our Method & \textbf{0.92} & \textbf{0.88} \\
    \bottomrule
  \end{tabular}
\end{table}

\section{Conclusion}

Summarize your work and discuss future directions.

%% Acknowledgments
\begin{acks}
This work was supported by [funding information].
\end{acks}

%% Bibliography
\bibliographystyle{ACM-Reference-Format}
\bibliography{references}

\end{document}
